\exercisesection{\S~4}

记号与假定同 §4 第 1、2、3 节。

\begin{enumerate}
\def\labelenumi{\arabic{enumi})}
\tightlist
\item
  取 \(\mathrm{G} = \mathbf{GL}_{n}(k), n \geq 0.\)
\end{enumerate}

\begin{enumerate}
\def\labelenumi{\alph{enumi})}
\item
  证明 \(r_{\mathrm{Ad}}^0 (g) = n\) 对一切 \(g \in \mathrm{G}\) 成立。
\item
  证明元素 \(g \in G\) 正则，当且仅当其特征多项式 \(\mathrm{P}_{g}(\mathrm{T}) = \det(\mathrm{T} - g)\) 可分；这等价于说 \(\mathrm{P}_{g}(\mathrm{T})\) 的判别式（《代数学》第四章 §1 第 10 节） \(\neq 0\) 。
\end{enumerate}

\begin{enumerate}
\def\labelenumi{\arabic{enumi})}
\setcounter{enumi}{1}
\item
  构造一个李群 G，使函数 \(r_{Ad}^{0}\) 不是常数。（取 g 交换且 \(\neq 0\) ，且 Ad 非平凡。）
\item
  设 \((\rho_{i})_{i\in I}\) 是 G 的可数族解析线性表示。证明 G 中对所有 \(\rho_{i}\) 都正则的元素构成 G 的稠密子集。构造一个例子说明可数性假定不可省去。
\item
  假设 \(k = \mathbf{C}\) 且 \(\mathbf{G}\) 连通。证明下列性质等价：
\end{enumerate}

\begin{enumerate}
\def\labelenumi{(\roman{enumi})}
\item
  G 幂零。
\item
  每个 \(\neq 1\) 的 \(G\) 中元素都正则。
\end{enumerate}

（先证明 (ii) 蕴涵

\begin{enumerate}
\def\labelenumi{(\roman{enumi})}
\setcounter{enumi}{1}
\tightlist
\item
  \(^{\prime}\) g 的每个 \(\neq 0\) 的元素都正则。
\end{enumerate}

再注意，若 \(g \neq 0\) ，则 g 含有元素 \(x \neq 0\) 使 ad x 幂零，参见 §3 习题 10。推出 g 幂零，从而得 (i)。

